\documentclass[letterpaper,10pt]{memoir}

% === AUTOR === (((
\author{\textit{Por Erick I. Rodríguez Juárez.}}
% )))

% === PAQUETES === (((
% \usepackage{makeidx}
% \usepackage{xltxtra}
\usepackage{amsfonts}
\usepackage{amsmath}
\usepackage{amssymb}
% \usepackage{fullpage}
\usepackage{tikz}
\usetikzlibrary{arrows.meta}
\usepackage{graphicx}
% )))

% === TIPOGRAFÍA === (((
% \setmainfont[
  % BoldFont       = bodonibi,
	% ItalicFont     = Century modern italic2.ttf,
	% BoldItalicFont = bodonibi,
	% SmallCapsFont  = lmromancaps10-regular.otf
% ]{Century_modern.ttf}
% )))

% === COMANDOS === (((
% \newcommand{\dis}{\displaystyle}
% \newcommand{\qed}{\hspace{0.5cm}\rule{0.16cm}{0.4cm}}
% \newcommand{\operator}[1]{\mathop{\vphantom{\sum}\mathchoice
% {\vcenter{\hbox{\huge $#1$}}}
% {\vcenter{\hbox{\Large $#1$}}}{#1}{#1}}\displaylimits}
% \newcommand{\suma}{\operator{\includegraphics[scale=0.09]{FOTOS/Sigma.png}}}
% \setlength{\parindent}{0mm}
% )))

% === ITALICA EN ENTORNO MATEMÁTICO === (((
% \DeclareSymbolFont{italics}{\encodingdefault}{\rmdefault}{m}{it}
% \DeclareSymbolFontAlphabet{\mathit}{italics}
% \ExplSyntaxOn
% \int_step_inline:nnnn { `A } { 1 } { `Z }
 % {  \exp_args:Nf \DeclareMathSymbol{\char_generate:nn{#1}{11}}{\mathalpha}{italics}{#1} }
% \int_step_inline:nnnn { `a } { 1 } { `z } {  \exp_args:Nf \DeclareMathSymbol{\char_generate:nn{#1}{11}}{\mathalpha}{italics}{#1}}
% \ExplSyntaxOff
% )))

\begin{document}

\titulo

\section*{Ecuaciones Diferenciales Homogéneas (12).} % (((
\begin{enumerate}
	\item \textbf{(Hacer 4 ejercicios)} Determine la solución general de las siguientes ecuaciones diferenciales de segundo orden.
		\begin{enumerate}
			\item \(4 y'' +y \,' =0\). \textbf{Solución.} \(y=c_1+c_2e^{-x/4}\).
			\item \(y'' -y \,' -6y=0\). \textbf{Solución.} \(y=c_1e^{3x} +c_2e^{-2x}\).
			\item \(y'' +8y \,' +16y=0\). \textbf{Solución.} \(y=c_1e^{-4x} +c_2xe^{-4x}\).
			\item \(12y '' -5y \,' -2y=0\). \textbf{Solución.} \(y=c_1e^{2x/3} +c_2e^{-x/4}\).
			\item \(y'' +9y=0\). \textbf{Solución.} \(y=c_1 \cos 3x+c_2 \sin 3x\).
			\item \(y'' -4y \,' +5y=0\). \textbf{Solución.} \(y=e^{2x} (c_1 \cos x+c_2 \sin x)\).
			\item \(3y'' +2y \,' +y=0\). \textbf{Solución.} \(y=e^{-x/3} \Bigg(c_1 \cos \dfrac{\sqrt{2}}{3} x+c_2 \sin \dfrac{\sqrt{2}}{3} x\Bigg)\).
		\end{enumerate}
	\item \textbf{(Hacer 4 ejercicios)} Determine la solución general de las siguientes ecuaciones diferenciales de orden superior
		\begin{enumerate}
			\item \(y''' -4y'' -5y \,' =0\). \textbf{Solución.} \(y=c_1+c_2e^{-x} +c_3e^{5x}\).
			\item \(y''' -5y'' +3y \,' +9y=0\). \textbf{Solución.} \(y=c_1e^{-x} +c_2e^{3x} +c_3xe^{3x}\).
			\item \(\dfrac{d^3u}{dt^3} + \dfrac{d^2u}{dt^2} 2u=0\). \textbf{Solución.} \(u=c_1e^t+e^{-t} (c_2 \cos t+c_3 \sin t)\).
			\item \(y''' +3y'' +3y \,' +y=0\). \textbf{Solución.} \(y=c_1e^{-x} +c_2xe^{-x} +c_3x^2e^{-x}\).
			\item \(y^{(4)} +y''' +y'' =0\). \textbf{Solución.} \(y=c_1+c_2x+e^{-x/2} \Big(c_3 \cos \dfrac{\sqrt{3}}{2} x+c_4 \sin \dfrac{\sqrt{3}}{2} x\Big)\).
			\item \(16 \dfrac{d^4y}{dt^4} +24 \dfrac{d^2 y}{dt^2} +9y=0\). \textbf{Solución.} \(y=c_1 \cos \dfrac{\sqrt{3}}{2} t+c_2 \sin \dfrac{\sqrt{3}}{2} t+c_3t \cos \dfrac{\sqrt{3}}{2} t+c_4 t \sin \dfrac{\sqrt{3}}{2} t\).
			\item \(\dfrac{d^5u}{dr^5} +5 \dfrac{d^4 u}{dr^4} -2 \dfrac{d^3u}{dr^3} -10 \dfrac{d^2u}{dr^2} + \dfrac{du}{dr} +5u=0\). \textbf{Solución.} \(u=c_1e^r+c_2re^r+c_3e^{-r} +c_4re^{-r} +c_5e^{-5r}\).
		\end{enumerate}
	\item \textbf{(Hacer 2 ejercicios)} Resuelve los siguientes problemas de valor inicial y esboza su grafica, haciendo énfasis en la interpretación geométrica de las condiciones iniciales.
		\begin{enumerate}
			\item \(y'' +16y=0\), \(y(0) =2\), \(y \,' (0) =-2\). \textbf{Solución.} \(y=2 \cos 4x- \dfrac{1}{2} \sin 4x\).
			\item \(\dfrac{d^2y}{dt^2} -4 \dfrac{dy}{dt} -5y=0\), \(y(1) =0\), \(y \,' (1) =2\). \textbf{Solución.} \(y=- \dfrac{1}{3} e^{-(t-1)} + \dfrac{1}{3} e^{5(t-1)}\).
			\item \(y'' +y \,' +2y=0\), \(y(0) =0\), \(y \,' (0) =0\). \textbf{Solución.} \(y=0\).
			\item \(y''' +12y'' +36y' =0\), \(y(0) =0\), \(y \,' (0) =0\), \(y'' (0) =-7\). \textbf{Solución.} \(y=- \dfrac{7}{36} + \dfrac{7}{36} e^{-6x} + \dfrac{7}{6} xe^{-6x}\).
		\end{enumerate}
	\item \textbf{(Hacer 2 ejercicios)} Resuelve las siguientes ecuaciones diferenciales, sujeta a las condiciones de frontera indicadas y esboza su gráfica, haciendo énfasis en la interpretación geométrica de las condiciones de frontera.
		\begin{enumerate}
			\item \(y'' -10y' +25y=0\), \(y(0) =1\), \(y(1) =0\). \textbf{Solución.} \(y=e^{5x} -xe^{5x}\).
			\item \(y'' +y=0\), \(y' (0) =0\), \(y' \Bigg(\dfrac{\pi}{2}\Bigg) =2\). \textbf{Solución.} \(y=-2 \cos x\).
		\end{enumerate}
\end{enumerate}
*Puedes usar GeoGebra \url{https://www.geogebra.org/m/KGWhcAqc} o WolframAlpha \url{https://www.wolframalpha.com/} para esbozar o verificar las graficas de las soluciones de PVI o PVF.
% )))

\end{document}
